\documentclass[11pt]{article}
% teacher_preamble.tex — shared preamble for the chalk-teaching TEACHER'S NOTES.
% These are standalone, instructor-facing lesson plans (article class), separate
% from the student lecture notes.  Compile each teaching/lecNN_teacher_notes.tex
% directly with pdfLaTeX.
\usepackage[T1]{fontenc}
\usepackage[utf8]{inputenc}
\usepackage[margin=1in]{geometry}
\usepackage{amsmath,amssymb,amsthm,mathtools}
\usepackage{braket}
\usepackage{bm}
\usepackage{graphicx}
\usepackage{enumitem}
\usepackage{xcolor}
\usepackage{tikz}
\usetikzlibrary{arrows.meta}
\usepackage{quantikz}
\usepackage[most]{tcolorbox}
\usepackage[colorlinks=true,linkcolor=blue,urlcolor=blue]{hyperref}

% ---- Notation macros (consistent with the course preamble.tex) ----
\newcommand{\R}{\mathbb{R}}
\newcommand{\C}{\mathbb{C}}
\newcommand{\N}{\mathbb{N}}
\newcommand{\Z}{\mathbb{Z}}
\newcommand{\E}{\mathbb{E}}
\newcommand{\Prob}{\mathbb{P}}
\newcommand{\Hilb}{\mathcal{H}}
\newcommand{\ii}{\mathrm{i}}
\newcommand{\dd}{\mathrm{d}}
\newcommand{\ee}{\mathrm{e}}
\DeclareMathOperator{\tr}{tr}
\DeclareMathOperator{\Tr}{Tr}
\DeclareMathOperator{\Var}{Var}
\DeclareMathOperator{\diag}{diag}
\newcommand{\norm}[1]{\left\lVert #1 \right\rVert}
\newcommand{\abs}[1]{\left\lvert #1 \right\rvert}
\newcommand{\bigo}{\mathcal{O}}
\newcommand{\id}{\mathbb{I}}
\newcommand{\dg}{^{\dagger}}
% quantikz also defines \proj (a circuit element we do not use); override it with our projector.
\let\proj\relax
\newcommand{\proj}[1]{\ket{#1}\!\bra{#1}}
\newcommand{\Xgate}{\mathsf{X}}
\newcommand{\Ygate}{\mathsf{Y}}
\newcommand{\Zgate}{\mathsf{Z}}
\newcommand{\Hgate}{\mathsf{H}}
\newcommand{\CNOT}{\mathrm{CNOT}}
\newcommand{\Uf}{U_f}

% ---- Theorem-like environments (numbered per section, article style) ----
\theoremstyle{definition}
\newtheorem{definition}{Definition}[section]
\newtheorem{example}[definition]{Example}
\theoremstyle{plain}
\newtheorem{theorem}[definition]{Theorem}
\newtheorem{lemma}[definition]{Lemma}
\newtheorem{proposition}[definition]{Proposition}
\theoremstyle{remark}
\newtheorem{remark}[definition]{Remark}

% ---- Teaching-specific structure ----
% Session block header with a time budget, e.g. \block{Qubits and the Bloch sphere}{45}
\newcommand{\block}[2]{%
  \bigskip\noindent
  \begin{tcolorbox}[colback=blue!6,colframe=blue!55!black,boxrule=0.6pt,arc=2pt,
    left=8pt,right=8pt,top=5pt,bottom=5pt]
  \textbf{\large #1}\hfill\textbf{[\,#2\,min]}
  \end{tcolorbox}}

% Board-plan cue: what to write/draw on the board.
\newtcolorbox{board}{colback=gray!8,colframe=gray!55!black,title={On the board},
  fonttitle=\bfseries\small,boxrule=0.5pt,arc=2pt,left=6pt,right=6pt,top=3pt,bottom=3pt}

% Graphic-demonstration cue (what to sketch).
\newcommand{\sketch}[1]{\par\smallskip\noindent{\color{teal}\textbf{Sketch:}} #1\par\smallskip}

% Optional / advanced material that may be skipped (but indicated).
\newtcolorbox{advanced}{colback=orange!7,colframe=orange!65!black,
  title={Advanced --- may be skipped (mention only)},
  fonttitle=\bfseries\small,boxrule=0.5pt,arc=2pt,left=6pt,right=6pt,top=3pt,bottom=3pt}

% Timing hint inline, e.g. \tmark{10} for a ~10-minute chunk.
\newcommand{\tmark}[1]{{\small\color{blue!55!black}(\,$\approx$#1\,min)}}

% In-class problem (numbered per section) and a teacher-facing answer.
\newtheorem{inclass}[definition]{In-class problem}
\newenvironment{answer}{\par\smallskip\noindent\textit{Answer.}\ }{\par\smallskip}

\title{\textbf{Lecture 3 --- Grover's Algorithm and Amplitude Amplification}\\[0.3em]\large Teacher's Notes (chalk, one 3$\times$45\,min meeting)}
\author{Course Instructor}
\date{}
\begin{document}
\maketitle

\noindent\textbf{How to use these notes.} This is a chalk-delivery lesson plan for
one three-hour meeting, split into three 45-minute blocks. Boxed \textsf{On the
board} panels give the exact equations and text to write; \textcolor{teal}{\textbf{Sketch}}
cues say what to draw; \tmark{n} markers pace each chunk. Work the derivations
\emph{live} at the board rather than reading them out. Each block has in-class
problems with teacher-facing answers---put the problem up, give students 3--6
minutes, then walk the solution. Orange \emph{Advanced} boxes are optional; mention
them only if time allows. The engine of the whole lecture is one picture ---
\emph{two reflections make a rotation} --- so keep the invariant-plane diagram up on
the side board once you draw it in Block 2.

\medskip
\noindent\textbf{Learning goals (by the end of the meeting a student can):}
\begin{itemize}[leftmargin=1.4em,nosep]
  \item State the bit and phase oracles, and derive phase kickback with a $\ket{-}$ ancilla.
  \item Run Deutsch, Deutsch--Jozsa, and Bernstein--Vazirani, explaining why one query suffices.
  \item Build Grover's two reflections (oracle $R_{x_0}$ and diffusion $R_{\psi_0}$) and prove $G$ rotates by $\theta$ in a 2-D plane.
  \item Compute the optimal iteration count $k^{\ast}\approx(\pi/4)\sqrt{N}$ (and $(\pi/4)\sqrt{N/M}$) and the success probability.
  \item State amplitude amplification ($\bigo(1/\sqrt p)$) and the $\Omega(\sqrt N)$ lower bound.
\end{itemize}

\medskip
\noindent\textbf{Prerequisites (Lectures 1--2).} Qubits as unit vectors, tensor
products, the standard gates $X,Z,H$, $H\ket0=\ket{+}$, $H\ket1=\ket{-}$,
computational-basis measurement (Born rule), and unitarity/reversibility. No prior
exposure to oracles assumed.

\medskip
\noindent\textbf{Materials.} Blackboard with room for a persistent ``cheat strip''
(oracle definitions, $H^{\otimes n}$ identity) on one side; \emph{colored chalk}
(white for algebra, one color for signs/phases created by the oracle, one for the
rotating state vector in the invariant plane).

% =====================================================================
\block{Block 1 --- Oracles, phase kickback, and the first speedups}{45}
% =====================================================================
\noindent\emph{Objective:} an oracle is a black-box unitary charged per query;
phase kickback turns a bit oracle into a sign; two Hadamard layers read the signs
out by interference---the template behind Deutsch, Deutsch--Jozsa, and
Bernstein--Vazirani.

\medskip
\noindent\textbf{Pacing:} the oracle/query model \tmark{8}; phase kickback (derive)
\tmark{9}; in-class problem 1 \tmark{5}; the Hadamard transform \& Deutsch
\tmark{9}; Deutsch--Jozsa \& Bernstein--Vazirani \tmark{9}; in-class problem 2
\tmark{5}.

\subsection*{1. The oracle / query model \tmark{8}}
Open on the cheat strip and leave it up all lecture: we do not see the code of a
Boolean $f$, only a reversible black box, and we count \emph{queries}, not gates.

\begin{board}
\textbf{Bit (XOR) oracle.} For $f:\{0,1\}^n\to\{0,1\}$, on input register $\ket x$
and target $\ket y$,
\[
  \Uf\ket{x}\ket{y}=\ket{x}\ket{y\oplus f(x)} .
\]
The $\oplus$ (not overwrite) makes it reversible: $\Uf^2=\id$, so $\Uf$ is unitary
and its own inverse. \textbf{Query complexity} $=$ number of calls to $\Uf$.
\textbf{Phase oracle} (what we really want): $\ket{x}\mapsto(-1)^{f(x)}\ket{x}$.
\end{board}

Stress the two-line moral: a classical query returns one pair $(x,f(x))$; a
quantum circuit can apply $\Uf$ to a \emph{superposition} of all inputs at once.
But raw parallelism is useless---measuring gives one random pair. The missing
ingredient is \emph{interference}, and phase kickback is how we create it.

\subsection*{2. Phase kickback --- derive it live \tmark{9}}
This is the single most important trick of the day. Put the target in $\ket{-}$ and
watch the sign appear on the input.

\begin{board}
\textbf{Phase kickback.} With $\ket{-}=\tfrac1{\sqrt2}(\ket0-\ket1)$,
\[
  \Uf\ket{x}\ket{-}
   =\tfrac1{\sqrt2}\big(\ket{x}\ket{f(x)}-\ket{x}\ket{1\oplus f(x)}\big)
   =(-1)^{f(x)}\ket{x}\ket{-}.
\]
\emph{Check the two cases:} $f(x)=0\Rightarrow\tfrac1{\sqrt2}(\ket0-\ket1)=+\ket{-}$;
$f(x)=1\Rightarrow\tfrac1{\sqrt2}(\ket1-\ket0)=-\ket{-}$. The target is unchanged
and unentangled, so it \emph{factors out}: on the input register $\Uf$ acts as the
phase oracle $\tilde U_f\ket x=(-1)^{f(x)}\ket x$.
\end{board}

\noindent The $\ket{-}$ ancilla that powers kickback ($X\ket{-}=-\ket{-}$, so the
target is a $-1$ eigenstate of the flip):
\begin{center}
\begin{quantikz}
\lstick{$\ket{x}$ (input)} & \gate[2]{\Uf} & \rstick{$(-1)^{f(x)}\ket{x}$}\qw \\
\lstick{$\ket1$}           & \gate{H}      & \qw
\end{quantikz}
\end{center}
The first $H$ prepares $\ket-$ from $\ket1$; the sign is kicked back onto the input.

\begin{inclass}[Compute a kickback action]\label{ic:l3-kickback}
Let $n=2$ and $f(x)=x_1\wedge x_2$ (AND), so $f=1$ only on $x=11$. Compute
$\Uf\big(\tfrac12\sum_{x}\ket x\big)\ket{-}$ and read off the resulting state of the
input register.
\end{inclass}
\begin{answer}
By kickback each branch picks up $(-1)^{f(x)}$, and the $\ket-$ factors out. Only
$x=11$ is marked, so the input register becomes
$\tfrac12(\ket{00}+\ket{01}+\ket{10}-\ket{11})$---a single sign flip on $\ket{11}$.
\checkmark
\end{answer}

\subsection*{3. The Hadamard transform and Deutsch \tmark{9}}
\begin{board}
\textbf{Hadamard identity.} For $x\in\{0,1\}^n$, with $x\cdot z=\bigoplus_i x_iz_i$,
\[
  H^{\otimes n}\ket{x}=\frac{1}{\sqrt{2^{\,n}}}\sum_{z\in\{0,1\}^n}(-1)^{x\cdot z}\ket{z},
  \qquad
  H^{\otimes n}\ket{0^n}=\frac{1}{\sqrt N}\sum_{x}\ket{x}\ (N=2^n).
\]
\textbf{Orthogonality:}\quad
$\dfrac1{2^n}\sum_{x}(-1)^{x\cdot w}=\delta_{w,0^n}$ (the interference is exact).
\end{board}

\begin{board}
\textbf{Deutsch ($n=1$).} Decide constant ($f(0)=f(1)$) vs.\ balanced
($f(0)\neq f(1)$) in \emph{one} query. Run $H$ on $\ket0$, query, $H$ again:
\[
  \ket{+}\ket{-}\xrightarrow{\Uf}
  \tfrac1{\sqrt2}\big((-1)^{f(0)}\ket0+(-1)^{f(1)}\ket1\big)\ket{-}
  \xrightarrow{H}
  \begin{cases}\pm\ket0 & \text{constant}\\ \pm\ket1 & \text{balanced.}\end{cases}
\]
Measure: outcome $0\Rightarrow$ constant, $1\Rightarrow$ balanced (classically needs
$2$ queries).
\end{board}

\subsection*{4. Deutsch--Jozsa and Bernstein--Vazirani \tmark{9}}
Same circuit, $n$ qubits: $H^{\otimes n}$, one query (target in $\ket-$),
$H^{\otimes n}$, measure. This template recurs all the way to Shor.

\begin{center}
\begin{quantikz}
\lstick{$\ket{0}^{\otimes n}$} & \gate{H^{\otimes n}} & \gate[2]{\Uf} & \gate{H^{\otimes n}} & \meter{} \\
\lstick{$\ket1$}               & \gate{H}             &               & \qw                  & \qw
\end{quantikz}
\end{center}

\begin{board}
\textbf{Deutsch--Jozsa} (promise: $f$ constant or balanced). After the query the
input is $\tfrac1{\sqrt N}\sum_x(-1)^{f(x)}\ket x$; a second $H^{\otimes n}$ gives
amplitude of $\ket{0^n}$
\[
  \alpha_{0^n}=\frac1N\sum_x(-1)^{f(x)}
  =\begin{cases}\pm1 & \text{constant}\\ 0 & \text{balanced.}\end{cases}
\]
So measuring $0^n$ certifies ``constant''; any nonzero outcome certifies ``balanced.''
\medskip

\textbf{Bernstein--Vazirani.} If $f_s(x)=s\cdot x$, the same circuit yields
$\ket{s}$ exactly: the post-query state is $H^{\otimes n}\ket s$, and a second
$H^{\otimes n}$ (its own inverse) returns $\ket s$. One query recovers all $n$ bits
(classically $n$).
\end{board}

\begin{inclass}[Determine the DJ outcome]\label{ic:l3-dj}
On $n=2$ take $f(x_1x_2)=x_1\oplus x_2$. Is $f$ constant or balanced? Trace the
Deutsch--Jozsa circuit and give the amplitude of $\ket{00}$ at the meter.
\end{inclass}
\begin{answer}
$f(00){=}0,f(01){=}1,f(10){=}1,f(11){=}0$: two $0$s and two $1$s, so
\emph{balanced}. Then $\alpha_{00}=\tfrac14\sum_x(-1)^{f(x)}
=\tfrac14(1-1-1+1)=0$, so $\ket{00}$ never occurs---correctly reporting
``balanced.'' (In fact $f=f_s$ with $s=11$, so the measured string is $11$.) \checkmark
\end{answer}

\begin{board}
\textbf{Who's who (say aloud).} \emph{David Deutsch} (Oxford): 1985 universal quantum
computer and the first quantum algorithm; a founder of the field. \emph{Richard Jozsa}
(Cambridge): the 1992 Deutsch--Jozsa generalization and a co-author of quantum
teleportation. \emph{Ethan Bernstein} \& \emph{Umesh Vazirani} (Berkeley): the 1993
\emph{Quantum complexity theory} paper that defined $\mathsf{BQP}$ and the first
$\mathsf{BPP}$--$\mathsf{BQP}$ oracle separation---of which the hidden-string algorithm
here is the seed.
\end{board}

\begin{advanced}
Bounded-error caveat: the dramatic ``$1$ vs.\ $2^{n-1}+1$'' Deutsch--Jozsa gap is a
\emph{deterministic} separation---a few random classical queries distinguish
constant from balanced with high confidence, so the randomized classical cost is
$\bigo(1)$. Bernstein--Vazirani's $1$-vs-$n$ separation survives even with bounded
error. Mention only.
\end{advanced}

% =====================================================================
\block{Block 2 --- Grover's algorithm}{45}
% =====================================================================
\noindent\emph{Objective:} unstructured search is two reflections---mark the
solution, then invert about the mean---and their product rotates the state toward
$\ket{x_0}$ by a fixed angle $\theta$ per step, so $\bigo(\sqrt N)$ steps suffice.

\medskip
\noindent\textbf{Pacing:} the problem \& phase oracle \tmark{7}; the two reflections
(board) \tmark{10}; the invariant plane \& rotation derivation \tmark{13}; in-class
problem 1 \tmark{5}; worked $N=8$ run \tmark{7}; in-class problem 2 \tmark{3}.

\subsection*{1. Unstructured search and the two reflections \tmark{17}}
\begin{board}
\textbf{Problem.} $N=2^n$ items; oracle marks a unique $x_0$ with $f(x_0)=1$, else
$0$. Classically $\Omega(N)$ queries (an adversary answers $0$ until one candidate
is left). Start state: uniform superposition
$\ket{\psi_0}=H^{\otimes n}\ket{0^n}=\tfrac1{\sqrt N}\sum_x\ket x$.
\end{board}

\begin{board}
\textbf{Two reflections.}
\[
  R_{x_0}=\id-2\proj{x_0}\quad(\text{phase-flip }\ket{x_0}),\qquad
  R_{\psi_0}=2\proj{\psi_0}-\id\quad(\text{fix }\ket{\psi_0},\text{ flip }\perp).
\]
Each is Hermitian, unitary, and an involution ($R^2=\id$): for $R=\id-2\proj u$,
$R\dg=R$ and $R^2=\id-4\proj u+4\proj u=\id$. \textbf{Grover iterate}
$G=R_{\psi_0}R_{x_0}$. Circuit form of diffusion:
\[
  R_{\psi_0}=H^{\otimes n}\bigl(2\proj{0^n}-\id\bigr)H^{\otimes n}.
\]
\emph{Inversion about the mean:} $R_{\psi_0}\sum_x a_x\ket x=\sum_x(2\bar a-a_x)\ket x$,
$\bar a=\tfrac1N\sum_x a_x$.
\end{board}

\noindent One Grover iterate as a circuit --- phase oracle $R_{x_0}$, then the
diffusion sandwich $H^{\otimes n}$, reflect-about-$\ket{0^n}$, $H^{\otimes n}$:
\begin{center}
\begin{quantikz}
\lstick{$\ket{\psi}$} & \gate{R_{x_0}} & \gate{H^{\otimes n}} & \gate{2\proj{0^n}-\id} & \gate{H^{\otimes n}} & \rstick{$G\ket{\psi}$}\qw
\end{quantikz}
\end{center}
Say aloud: $R_{x_0}$ is the one oracle call; everything else is cheap gates. In real
uses $x_0$ is \emph{not} hardcoded---the oracle tests a predicate ``is this a
solution?'' and marks whoever passes.

\begin{board}
\textbf{Compiling $R_{x_0}$ from a verifier} (never from $x_0$). For ``find $x$ with
$g(x)=t$'': (1) \emph{compute} $\ket{x}\ket{0^m}\mapsto\ket{x}\ket{g(x)}$ with a
reversible circuit $U_g$; (2) \emph{mark} the known target $t$ ($\Xgate$-conjugated
multi-controlled $\Zgate$ on the output register); (3) \emph{uncompute} $U_g^{\dagger}$.
Net: $\ket{x}\mapsto(-1)^{[\,g(x)=t\,]}\ket{x}=R_{x_0}$, with $x_0=g^{-1}(t)$ never
found. \emph{Recognize, don't locate.}
\end{board}
\noindent\emph{Example oracles} (say aloud): key search $g(k)=\mathrm{Enc}_k(P_0)$,
$t$ the ciphertext (AES-128: $2^{128}\!\to\!2^{64}$, why keys double);
hash preimage $g=H$; satisfiability $g=\varphi,\ t=1$; unstructured table
$g(x)=L[x]$. Grover wins only when $g$ is easy to \emph{check} but has no structure
to exploit (else Shor/BV/convex methods beat it).

\begin{board}
\textbf{The invariant plane.} Let $\ket{\psi_0^\perp}=\tfrac1{\sqrt{N-1}}\sum_{x\neq x_0}\ket x$.
Then
\[
  \ket{\psi_0}=\tfrac1{\sqrt N}\ket{x_0}+\tfrac{\sqrt{N-1}}{\sqrt N}\ket{\psi_0^\perp}
  =\sin\tfrac\theta2\ket{x_0}+\cos\tfrac\theta2\ket{\psi_0^\perp},
  \quad \sin\tfrac\theta2=\tfrac1{\sqrt N}.
\]
$\mathcal P=\operatorname{span}\{\ket{x_0},\ket{\psi_0^\perp}\}$. In the ordered basis
$\{\ket{\psi_0^\perp},\ket{x_0}\}$:
\[
  R_{x_0}=\begin{pmatrix}1&0\\0&-1\end{pmatrix},\qquad
  R_{\psi_0}=\begin{pmatrix}\cos\theta&\sin\theta\\ \sin\theta&-\cos\theta\end{pmatrix}
  \ (\text{reflect about }\ket{\psi_0}\text{ at angle }\tfrac\theta2).
\]
\[
  G=R_{\psi_0}R_{x_0}
   =\begin{pmatrix}\cos\theta&-\sin\theta\\ \sin\theta&\cos\theta\end{pmatrix}
   \ =\ \textbf{rotation by }\theta.
\]
\end{board}

Derive the rotation angle live: $R_{x_0}$ is a reflection about the
$\ket{\psi_0^\perp}$ axis; $R_{\psi_0}$ is a reflection about $\ket{\psi_0}$, which
sits at angle $\theta/2$ from $\ket{\psi_0^\perp}$. Product of two reflections $=$
rotation by \emph{twice} the angle between their mirror lines $=2\cdot\tfrac\theta2=\theta$.

\sketch{The invariant plane. Horizontal axis $\ket{\psi_0^\perp}$ (``bad''),
vertical axis $\ket{x_0}$ (``good''). Draw $\ket{\psi_0}$ close to horizontal at
angle $\theta/2$ above the $\ket{\psi_0^\perp}$ axis (since $\sin\tfrac\theta2=1/\sqrt N$
is small). Draw a curved arrow lifting the state by $\theta$ each iterate; after $k$
iterates it sits at angle $(2k{+}1)\tfrac\theta2$ from horizontal. \emph{Say aloud:}
each step tips the vector toward $\ket{x_0}$; we stop when it is nearly vertical.}

\begin{board}
\textbf{State after $k$ iterations.}
\[
  G^k\ket{\psi_0}=\sin\!\tfrac{(2k+1)\theta}{2}\,\ket{x_0}+\cos\!\tfrac{(2k+1)\theta}{2}\,\ket{\psi_0^\perp},
  \qquad
  p_k=\sin^2\!\tfrac{(2k+1)\theta}{2}.
\]
Stop when $(2k{+}1)\theta/2\approx\pi/2$: $k^\ast=\lfloor(\pi/4)\sqrt N\rfloor$,
$p_{k^\ast}\ge1-1/N$. So $\bigo(\sqrt N)$ queries (quadratic speedup over $\bigo(N)$).
\end{board}

\begin{inclass}[Count Grover iterations]\label{ic:l3-count}
For $N=2^{20}$, how many iterations $k^\ast$ maximize the success probability? How
does this compare to classical search?
\end{inclass}
\begin{answer}
$k^\ast=\lfloor(\pi/4)\sqrt{2^{20}}\rfloor=\lfloor(\pi/4)\cdot1024\rfloor
=\lfloor 804.2\rfloor=804$ iterations (one oracle call each). Classical search needs
$\bigo(N)=\bigo(2^{20})\approx10^6$ queries: about $2^{10}\approx1000\times$ more. \checkmark
\end{answer}

\subsection*{2. A fully worked run: $N=8$, $x_0=3$ \tmark{7}}
Track two numbers: the common unmarked amplitude and the marked amplitude
$g$. Oracle flips the sign of $g$; diffusion sends every $a_x\mapsto2\bar a-a_x$.
Here $\sin(\theta/2)=1/\sqrt8\Rightarrow\theta\approx41.4^\circ$, and
$k^\ast=\lfloor(\pi/4)\sqrt8\rfloor=\lfloor2.22\rfloor=2$.

\begin{board}
\[
\begin{array}{lccc}
\text{stage} & \text{unmarked} & \text{marked }\ket3 & p=\abs g^2\\[2pt]
\text{start }\ket{\psi_0} & 1/\sqrt8 & 1/\sqrt8 & 1/8\\
\text{after }R_{x_0}\ (\text{q1}) & 1/\sqrt8 & -1/\sqrt8 & -\\
\text{after }R_{\psi_0}\ (\text{it.1}) & 1/(2\sqrt8) & 5/(2\sqrt8) & 25/32\approx0.781\\
\text{after }R_{x_0}\ (\text{q2}) & 1/(2\sqrt8) & -5/(2\sqrt8) & -\\
\text{after }R_{\psi_0}\ (\text{it.2}) & -1/(4\sqrt8) & 11/(4\sqrt8) & 121/128\approx0.945
\end{array}
\]
\end{board}

Do iteration~1 live: after the oracle $\bar a=(7\cdot\tfrac1{\sqrt8}-\tfrac1{\sqrt8})/8
=\tfrac{3}{4\sqrt8}$; inversion gives unmarked $2\bar a-\tfrac1{\sqrt8}=\tfrac1{2\sqrt8}$,
marked $2\bar a+\tfrac1{\sqrt8}=\tfrac5{2\sqrt8}$, so $p_1=25/32$. The negative
unmarked sign after iteration~2 means we have rotated just \emph{past} $\ket{x_0}$: a
third iterate would \emph{lose} probability (over-rotation).

\begin{inclass}[Grover on $N=4$ is exact]\label{ic:l3-n4}
For $N=4$, $M=1$, compute $\theta$ and show that a \emph{single} Grover iterate
reaches $\ket{x_0}$ with certainty.
\end{inclass}
\begin{answer}
$\sin(\theta/2)=1/\sqrt4=\tfrac12\Rightarrow\theta/2=\pi/6$, $\theta=\pi/3$. After
$k=1$ the angle is $(2\cdot1+1)\theta/2=\tfrac{3}{2}\cdot\tfrac\pi3=\tfrac\pi2$, so
$G\ket{\psi_0}=\sin\tfrac\pi2\ket{x_0}+\cos\tfrac\pi2\ket{\psi_0^\perp}=\ket{x_0}$:
$p_1=1$, deterministic. \checkmark
\end{answer}

% =====================================================================
\block{Block 3 --- Amplitude amplification, optimality, and applications}{45}
% =====================================================================
\noindent\emph{Objective:} $M$ solutions speed the search to $(\pi/4)\sqrt{N/M}$;
the two-reflection trick generalizes to \emph{any} state preparation (amplitude
amplification, $\bigo(1/\sqrt p)$); and a hybrid argument shows $\Omega(\sqrt N)$ is
optimal.

\medskip
\noindent\textbf{Pacing:} multiple solutions \tmark{9}; in-class problem 1
\tmark{5}; amplitude amplification (general trick) \tmark{11}; in-class problem 2
\tmark{5}; the $\Omega(\sqrt N)$ lower bound \tmark{8}; applications \& wrap-up
\tmark{7}.

\subsection*{1. Multiple solutions \tmark{9}}
\begin{board}
\textbf{$M$ marked items.} Good/bad states
$\ket{\mathrm{good}}=\tfrac1{\sqrt M}\sum_{x\in S}\ket x$,
$\ket{\mathrm{bad}}=\tfrac1{\sqrt{N-M}}\sum_{x\notin S}\ket x$, so
$\ket{\psi_0}=\sqrt{M/N}\,\ket{\mathrm{good}}+\sqrt{(N-M)/N}\,\ket{\mathrm{bad}}$.
Same rotation with
\[
  \sin\tfrac\theta2=\sqrt{\tfrac MN},\qquad
  k^\ast=\Big\lfloor\tfrac\pi4\sqrt{\tfrac NM}\Big\rfloor=\bigo\!\big(\sqrt{N/M}\big).
\]
More solutions $\Rightarrow$ faster search.
\end{board}

Warn (echo Block 2): if you \emph{overestimate} how many iterations to run, the
state rotates past $\ket{\mathrm{good}}$ and $p_k$ falls---always stop at $k^\ast$.

\begin{inclass}[Iterations with $M$ solutions]\label{ic:l3-M}
For $N=1024$ with $M=4$ marked items, find $\sin(\theta/2)$ and $k^\ast$, and
compare with the single-solution ($M=1$) count for the same $N$.
\end{inclass}
\begin{answer}
$\sin(\theta/2)=\sqrt{4/1024}=\sqrt{1/256}=1/16$. Then
$k^\ast=\lfloor(\pi/4)\sqrt{1024/4}\rfloor=\lfloor(\pi/4)\sqrt{256}\rfloor
=\lfloor(\pi/4)\cdot16\rfloor=\lfloor12.57\rfloor=12$. For $M=1$:
$k^\ast=\lfloor(\pi/4)\sqrt{1024}\rfloor=\lfloor(\pi/4)\cdot32\rfloor
=\lfloor25.13\rfloor=25$. Four solutions halve the count ($\sqrt M=2\times$ fewer). \checkmark
\end{answer}

\subsection*{2. Amplitude amplification --- the general trick \tmark{11}}
Grover is one instance of a reflection recipe that works for \emph{any} preparation
$\mathcal A$. Replace $H^{\otimes n}$ by $\mathcal A$ and $\ket{x_0}$ by the good
subspace.

\begin{board}
\textbf{Amplitude amplification.} Suppose
$\mathcal A\ket0=\sqrt{p}\,\ket{\mathrm{good}}+\sqrt{1-p}\,\ket{\mathrm{bad}}$ and a
projector $P$ recognizes the good part. Set
\[
  S_P=\id-2P,\qquad S_0=\mathcal A(2\proj0-\id)\mathcal A\dg,\qquad Q=S_0S_P.
\]
$Q$ rotates by $\theta$ with $\sin(\theta/2)=\sqrt p$ in
$\operatorname{span}\{\ket{\mathrm{good}},\ket{\mathrm{bad}}\}$; after
\[
  k^\ast=\Big\lfloor\tfrac{\pi}{4\arcsin\sqrt p}\Big\rfloor=\bigo\!\big(1/\sqrt p\big)
\]
iterations we get $\ket{\mathrm{good}}$ with probability $\Omega(1)$.
\end{board}

Say aloud: the proof is \emph{verbatim} the Grover proof --- $S_P$ reflects about
the bad axis, $S_0=\mathcal A(2\proj0-\id)\mathcal A\dg$ reflects about the axis
$\mathcal A\ket0$, and two reflections make a rotation. The payoff: classical
prepare-and-check needs $\bigo(1/p)$ tries to see the good outcome; amplitude
amplification needs $\bigo(1/\sqrt p)$---a quadratic win on \emph{any} subroutine,
not just search. Grover $=$ the case $\mathcal A=H^{\otimes n}$, $p=M/N$.

\begin{inclass}[Amplitude-amplification iteration count]\label{ic:l3-aa}
A preparation $\mathcal A$ succeeds with probability $p=1/100$ on a single shot. How
many amplitude-amplification iterations $k^\ast$ are needed, and how does that
compare with classical repetition?
\end{inclass}
\begin{answer}
$\sin(\theta/2)=\sqrt p=1/10$, so $\theta/2=\arcsin(0.1)=0.10017$ rad and
$k^\ast=\lfloor(\pi/4)/0.10017\rfloor=\lfloor7.84\rfloor=7$ iterations. Classical
prepare-and-check needs $\bigo(1/p)=\bigo(100)$ shots to expect one success: about
$1/\sqrt p=10\times$ fewer with amplification. \checkmark
\end{answer}

\subsection*{3. The $\Omega(\sqrt N)$ lower bound \tmark{8}}
State it, sketch the idea; keep the argument statement-level.

\begin{board}
\textbf{Lower bound.} Any quantum algorithm that finds $x_0$ (or even decides
whether \emph{a} marked item exists) with constant success probability needs
$\Omega(\sqrt N)$ queries. Combined with $k^\ast=\bigo(\sqrt N)$: Grover is optimal,
$\Theta(\sqrt N)$.
\end{board}

\begin{advanced}
\textbf{Hybrid / progress argument (idea only).} Let $\ket{\phi_k}$ be the state
after $k$ queries with the empty oracle $f\equiv0$, and $\ket{\phi_k^{x_0}}$ with
$x_0$ marked. Track
$D_k=\sum_{x_0}\norm{\ket{\phi_k^{x_0}}-\ket{\phi_k}}^2$. One query changes each
term by a bounded amount, giving $D_k\le 4k^2$. But to \emph{identify} the marked
item the branches must be nearly orthogonal, forcing $D_k=\Omega(N)$. Hence
$\Omega(N)\le 4k^2\Rightarrow k=\Omega(\sqrt N)$. A super-polynomial speedup would
require exploitable structure (Shor, phase estimation)---not available to an
unstructured oracle. Mention only.
\end{advanced}

\begin{inclass}[Why not run longer?]\label{ic:l3-overrotate}
Using $p_k=\sin^2((2k+1)\theta/2)$, explain what happens to the success probability
if one runs $2k^\ast$ iterations instead of $k^\ast$ (take $N=2^{20}$).
\end{inclass}
\begin{answer}
At $k^\ast$ the angle $(2k^\ast{+}1)\theta/2\approx\pi/2$ and $p\approx1$. Doubling
to $2k^\ast$ roughly doubles the angle to $\approx\pi$, where
$\sin^2(\pi)\approx0$: the state has rotated \emph{past} $\ket{x_0}$ and almost all
the way back to $\ket{\psi_0^\perp}$, so success collapses to nearly $0$. Grover is
periodic---more is \emph{not} better; stop at $k^\ast$. \checkmark
\end{answer}

\subsection*{4. Applications and wrap-up \tmark{7}}
Close by tying the trick to real problems.
\begin{itemize}[leftmargin=1.5em,itemsep=1pt,topsep=2pt]
  \item \emph{Unstructured search / database lookup:} $\bigo(\sqrt N)$ vs.\ $\bigo(N)$.
  \item \emph{NP-style brute force:} package ``is this assignment satisfying?'' as
        the oracle; e.g.\ a SAT instance on $n=30$ vars with $M=2^{10}$ solutions
        needs $\approx(\pi/4)\sqrt{2^{20}}\approx804$ queries vs.\ $\approx2^{20}$
        classically.
  \item \emph{Quantum counting:} apply phase estimation (next lecture) to the Grover
        operator $G$ to \emph{estimate} $M$ from the rotation angle $\theta$.
  \item \emph{Amplitude amplification as glue:} a quadratic boost wrapped around any
        randomized quantum subroutine with success probability $p$.
\end{itemize}
\noindent\textbf{One-sentence takeaway (say aloud):} operate on \emph{amplitudes},
not probabilities---the angle grows linearly in $k$, probability is its square, so a
target hit in $\bigo(N)$ probability-steps is hit in $\bigo(\sqrt N)$ amplitude-steps.

\begin{advanced}
\textbf{Bridge to Lecture 4 --- phase estimation.} The Grover iterate $G$ has
eigenvalues $\ee^{\pm\ii\theta}$ on the invariant plane; measuring that phase is
exactly quantum phase estimation, which turns Grover into \emph{counting} (estimate
$M$) and underlies the QFT-based algorithms next time. Mention only.
\end{advanced}

\end{document}
